% ============================================================================
%  chapter17.tex
%  Network Programming: Sockets and the Scalable Server
% ============================================================================

\chapter{Network Programming}

\section{Beyond One Machine}

Every channel in this book so far stayed inside one machine: pipes
between processes on the same kernel, files on the same disk, signals
to the same process tree. The network is the last boundary, and the
most important one. It lets two processes on different machines---on
different continents---exchange bytes, using a stack of protocols that
the operating system provides and mostly manages for you.

The astonishing fact that makes networking feel like the chapters
before it is this: \emph{a network connection looks like a file
descriptor}. You have already written the hard parts. After the setup,
reading from a socket is \code{read}, writing is \code{write}, and the
end of the conversation is the same EOF you met in Chapter~7. The new
machinery is only the setup: naming the machine, naming the program,
and agreeing on the rules of the conversation.

\begin{verbatim}
   The whole book, in one arrow:

     our process  --read/write-->  a FILE DESCRIPTOR
                                      |
                        +-------------+-------------+
                        |             |             |
                      a file       a pipe       a SOCKET
                      (ch.5)      (ch.7)    (this chapter)
                                              |
                                        another machine,
                                        thousands of km away
\end{verbatim}

\begin{osconcept}
The network is the operating system's biggest abstraction: the kernel
owns the network hardware, implements the protocol stacks, and
presents every connection as a descriptor your program reads and
writes exactly like a file. Only the OS can do this, because only the
OS can multiplex thousands of connections onto one network card and
run the intricate protocol machinery that makes bytes survive the
unreliable network between you and your peer.
\end{osconcept}

\section{The Three Names of a Network Conversation}

To reach a program on another machine, your operating system needs
three facts, and each is a different kind of name:

\begin{description}
  \item[The IP address] names the \emph{machine}. \code{192.168.1.10}
    is a 32-bit number (IPv4) written as four bytes. It is the
    network-layer name: ``that machine, out there.''
  \item[The port] names the \emph{program} on that machine. A machine
    runs many network programs at once; the port is how a connection
    chooses one of them. \code{80} is the conventional web port,
    \code{53} DNS, and your own servers use whatever you pick.
  \item[The protocol] names the \emph{rules of conversation}. The two
    that matter are TCP and UDP: TCP for reliable streams, UDP for
    fast datagrams. Both travel on top of the same IP addressing.
\end{description}

\begin{verbatim}
   A conversation needs all three:

       192.168.1.10 : 9000     over  TCP
        ^^^^^^^^^^^   ^^^^           ^^^
        machine      program      the rules

   The socket is exactly this triple. Create one, and you have asked
   the OS for "a network endpoint on this machine, using this
   protocol."
\end{verbatim}

\begin{osconcept}
Ports are an OS-level namespace, exactly like the process tree and the
filesystem hierarchy: the kernel keeps a table of ``which socket owns
which port'' and enforces that only one program is bound to a port at
a time. Without the kernel's arbitration, two programs could both
claim port 9000 and incoming packets would have nowhere unambiguous to
go. Naming and arbitration of ports is kernel work.
\end{osconcept}

\section{Sockets and Byte Order}

A \emph{socket} is an endpoint for network communication---the thing
your program holds as a descriptor. Creating one and giving it a name
and a port is the setup dance, and it involves one concept you have
not seen since Chapter~16: \emph{byte order}.

The network's wire format is defined as \emph{big-endian}: the most
significant byte first, regardless of what your CPU does. On a
little-endian machine (our x86-64), that means integers written into
socket addresses must be \emph{converted} before the kernel will read
them correctly. The helpers are \code{htons} (host to network, short)
and \code{htonl} (host to network, long). You met the problem in
Chapter~16; the network is where it is unavoidable.

\begin{lstlisting}[caption={Creating a socket address in network byte order.}]
#include <stdio.h>
#include <stdlib.h>
#include <string.h>
#include <unistd.h>
#include <arpa/inet.h>
#include <sys/socket.h>

#define PORT 9000

int main(void) {
    int fd;
    struct sockaddr_in addr;

    /* 1. Create the socket: IPv4 + TCP (SOCK_STREAM). */
    fd = socket(AF_INET, SOCK_STREAM, 0);
    if (fd < 0) {
        perror("socket");
        return 1;
    }

    /* 2. Build the address: family, port, interface. */
    memset(&addr, 0, sizeof(addr));
    addr.sin_family = AF_INET;                 /* IPv4 */
    addr.sin_port = htons(PORT);               /* port in NETWORK order */
    addr.sin_addr.s_addr = htonl(INADDR_ANY);  /* any interface */

    printf("socket created, fd = %d\n", fd);
    printf("port in network byte order: %d\n", addr.sin_port);
    close(fd);
    return 0;
}
\end{lstlisting}

\begin{verbatim}
   The socket address struct, drawn:

     struct sockaddr_in {
         short   sin_family;     AF_INET (IPv4) or AF_INET6 (IPv6)
         unsigned short sin_port;  the port, in BIG-ENDIAN byte order
         struct in_addr sin_addr;  the IP address, big-endian too
         ...                        (padding to the generic sockaddr)
     };

   - htons(9000) on a little-endian machine stores the bytes reversed,
     so that the kernel (and the wire) sees 9000 correctly.
   - htonl(INADDR_ANY) means "accept connections on ALL interfaces."
   - inet_pton(AF_INET, "127.0.0.1", &addr.sin_addr) converts a
     dotted-quad string into the 32-bit address (the client's job).
\end{verbatim}

\begin{ruleofthumb}[The wire is big-endian; your struct is not]
Whatever your CPU does, the network format is fixed: most significant
byte first. Every integer that leaves your machine---ports, addresses,
and any data you serialize---must be converted to network byte order,
or the machine on the other end will read your numbers backwards. This
is Chapter~16's wire-format discipline with a byte order the protocol,
not you, chooses.
\end{ruleofthumb}

\section{A TCP Server, Step by Step}

TCP is the \emph{connection-oriented} protocol: before any data, the
two sides perform a handshake, and after the data, they close. The
server's setup has four calls, and each has one job:

\begin{verbatim}
   The server lifecycle:

     socket()   create the endpoint              (a descriptor)
     bind()     give it a name: IP + port        ("I live here")
     listen()   start accepting; the kernel       ("queue my calls")
                queues incoming connections
     accept()   take one connection off the       (returns a NEW
                queue; returns a NEW socket        socket for THIS
                for that one client                client)

   accept() is the confusing step, and it is the heart of the model:
   the LISTENING socket is not for talking -- it is for meeting.
   Each accepted connection is a separate socket, with its own buffer
   and its own conversation.
\end{verbatim}

\begin{lstlisting}[caption={tcp\_server.c: a complete one-client echo server.}]
#include <stdio.h>
#include <stdlib.h>
#include <string.h>
#include <unistd.h>
#include <arpa/inet.h>
#include <sys/socket.h>

#define PORT 9000

int main(void) {
    int listen_fd;
    int conn_fd;
    struct sockaddr_in addr;
    char buf[1024];
    ssize_t got;

    listen_fd = socket(AF_INET, SOCK_STREAM, 0);
    if (listen_fd < 0) {
        perror("socket");
        return 1;
    }

    memset(&addr, 0, sizeof(addr));
    addr.sin_family = AF_INET;
    addr.sin_port = htons(PORT);
    addr.sin_addr.s_addr = htonl(INADDR_ANY);

    if (bind(listen_fd, (struct sockaddr *) &addr, sizeof(addr)) < 0) {
        perror("bind");
        return 1;
    }

    if (listen(listen_fd, 16) < 0) {        /* backlog: 16 queued */
        perror("listen");
        return 1;
    }

    printf("listening on port %d\n", PORT);

    conn_fd = accept(listen_fd, NULL, NULL);  /* blocks until a client */
    if (conn_fd < 0) {
        perror("accept");
        return 1;
    }

    printf("client connected\n");

    /* Echo whatever the client sends until it closes. */
    while (1) {
        got = read(conn_fd, buf, sizeof(buf));
        if (got <= 0) {
            break;                  /* EOF: the client closed */
        }
        write(conn_fd, buf, (size_t) got);
    }

    close(conn_fd);
    close(listen_fd);
    return 0;
}
\end{lstlisting}

The \emph{listen backlog} is the kernel's welcome mat: it tells the
kernel how many pending connections it may queue while you are busy
with others. If the queue is full, new connection attempts are simply
delayed or refused---the kernel, not your program, holds the door.

\begin{verbatim}
   Why the backlog exists, drawn:

        kernel's listen queue (capacity 16)
        +------+------+------+...
        | conn | conn | ...         <- the kernel holds these WHILE
        +------+------+------+...      your program handles others
        ^
     clients connecting

     your program:   accept()  takes ONE from the queue
                     read/write over it
                     accept()  takes the next...
\end{verbatim}

\begin{osconcept}
The listening socket, the backlog queue, and the handshake are all
kernel state. Your \code{listen} call only configures the kernel's
welcoming; the kernel itself accepts the TCP handshake, queues the
connection, and hands it to your \code{accept}. This is why a server
can be slow to \code{accept} without dropping connections---the OS is
holding the door. Only the kernel can do this, because only the kernel
owns the protocol stack.
\end{osconcept}

\section{A TCP Client: The Three-Way Handshake}

The client's setup is two calls: \code{socket} to create the endpoint,
and \code{connect} to reach the server. The magic happens inside
\code{connect}: it drives the TCP \emph{three-way handshake} with the
kernel's help, so that when \code{connect} returns, both sides agree a
connection exists and data can flow.

\begin{lstlisting}[caption={tcp\_client.c: send a message, print the reply.}]
#include <stdio.h>
#include <stdlib.h>
#include <string.h>
#include <unistd.h>
#include <arpa/inet.h>
#include <sys/socket.h>

#define PORT 9000

int main(int argc, char *argv[]) {
    int fd;
    struct sockaddr_in addr;
    char buf[1024];
    ssize_t got;

    if (argc != 2) {
        fprintf(stderr, "usage: %s message\n", argv[0]);
        return 1;
    }

    fd = socket(AF_INET, SOCK_STREAM, 0);
    if (fd < 0) {
        perror("socket");
        return 1;
    }

    memset(&addr, 0, sizeof(addr));
    addr.sin_family = AF_INET;
    addr.sin_port = htons(PORT);
    if (inet_pton(AF_INET, "127.0.0.1", &addr.sin_addr) <= 0) {
        perror("inet_pton");
        return 1;
    }

    /* connect performs the three-way handshake. */
    if (connect(fd, (struct sockaddr *) &addr, sizeof(addr)) < 0) {
        perror("connect");
        return 1;
    }

    write(fd, argv[1], strlen(argv[1]));

    got = read(fd, buf, sizeof(buf));      /* blocks until the reply */
    if (got > 0) {
        printf("server echoed: %.*s\n", (int) got, buf);
    }

    close(fd);
    return 0;
}
\end{lstlisting}

\begin{verbatim}
   The three-way handshake, drawn (the SYN packets):

     server  (LISTEN)                 client
        |                              socket()
        |<---------- SYN ------------   connect(): "can we talk?"
        |---------- SYN + ACK ------->  "yes, and my handshake"
        |<---------- ACK -------------  "agreed"
        |                              connect() returns: ESTABLISHED
        |<--------- data -------------  now the read/write phase
        |--------- data ------------->

   Why THREE messages? Because each side must learn that the other is
   there AND willing: the SYN asks, the SYN+ACK answers and asks back,
   and the ACK closes the loop. A two-message version could not tell a
   dead peer from a busy one.

   Closing is a two-way FIN exchange: each side sends FIN when it is
   done writing, and the other sees EOF (read returns 0). Half-closing
   is legal: one side can finish sending while still reading.
\end{verbatim}

\subsection{How a Connection Works in Bytes}

The handshake above is real packets, and packets are real bytes. Every
TCP message---handshake or data---is an IP packet whose payload is a
\emph{TCP segment}: a header followed by the user's data. The header
is a wire format, exactly like the record format of Chapter~16, and it
is what carries the connection's state. Here is the header, byte by
byte:

\begin{verbatim}
   The TCP header (minimal form), drawn like a wire format:

     bytes    field          meaning
     ------   -------------  -----------------------------------------
     0-1      src port       where THIS segment comes from
     2-3      dst port       where it is going
     4-7      sequence       the byte's position in the outgoing stream
     8-11     ack number     the next byte expected from the peer
     12       data offset    header length in 32-bit words (4 bits)
     13       flags          (URG ACK PSH RST SYN FIN, among others)
     14-15    window         flow control: how many bytes the receiver
                             is willing to accept
     16-17    checksum       integrity check over header and payload
     18-19    urgent pointer (meaningful only when URG is set)
     20+      payload        the user's data bytes (for data segments)

   One segment is therefore: 20 bytes of protocol + the data.
   The connection itself has NO bytes on the wire -- it IS the
   sequence numbers and flags shared across segments.
\end{verbatim}

The two numbers that make the byte stream work are the
\emph{sequence} and \emph{acknowledgment} fields. Every byte of data
is numbered: the sender stamps each segment with the sequence number
of its first byte, and the receiver acknowledges by sending back
``I expect byte N next.'' This one mechanism does three jobs at once:
ordering (the receiver sorts out-of-order segments by their numbers),
retransmission (an unacknowledged byte is resent), and the byte
stream your \code{read} sees.

\begin{verbatim}
   The three-way handshake in bytes:

     client                            server
        |  SYN, seq=1000                |
        |------------------------------>|   "hello, start at byte 1000"
        |  SYN+ACK, seq=7000, ack=1001  |
        |<------------------------------|   "hello, I start at 7000;
        |                               |    I expect your byte 1001"
        |  ACK, seq=1001, ack=7001      |
        |------------------------------>|   "agreed; I'm sending 1001"
        |                               |
        |  data "ada", seq=1001         |   <-- the payload bytes
        |------------------------------>|       (sequence 1001 = first
        |                               |        byte of "ada")
        |  ACK, ack=1004                |
        |<------------------------------|   "got bytes 1001..1003;
        |                               |    next I expect 1004"

   Read it as a wire format conversation: each arrow is a real
   packet, each packet carries sequence numbers that let the kernel
   on each side turn an arbitrary stream of segments into the ordered,
   gap-free byte stream your read() returns.
\end{verbatim}

The lesson, tying the whole book together: a TCP connection is a
\emph{protocol}, and a protocol is bytes with rules---the same idea as
the MiniDB wire format of Chapter~16, scaled to two machines. The
kernel implements the rules; your program sees only the resulting byte
stream.

\begin{osconcept}
The byte-level machinery---sequence numbers, acknowledgments,
retransmission, flow control---is the kernel's own wire-format
implementation. Your program never sees segments or sequence numbers;
the kernel consumes the wire protocol on your behalf and delivers a
clean byte stream. This is the extreme of the book's theme: the OS
turns raw, unreliable bytes on a wire into the abstractions your code
was built to use.
\end{osconcept}

\begin{osconcept}
The TCP handshake and connection state machine live entirely in the
kernel. Your \code{connect} is one syscall; behind it the kernel
retransmits lost SYN packets, times out dead peers, and tracks the
connection's state (SYN-SENT, ESTABLISHED, FIN-WAIT, and so on)
throughout. Reliability on an unreliable network is exactly the kind
of tedious, time-sensitive machinery no application could duplicate--
and exactly what an operating system is for.
\end{osconcept}

\section{Reading and Writing on a Socket}

After the setup, a connected socket is a file descriptor, and every
lesson of Chapters~4 and~6 applies unchanged:

\begin{verbatim}
   The socket is a byte stream, exactly like a pipe:

     - read() may return FEWER bytes than requested: loop or accept
       short reads.
     - EOF (read returns 0) means the peer closed: its FIN arrived.
     - write() may write fewer bytes than asked: loop.
     - the data has NO message boundaries. "ab" then "cd" may arrive
       as "abcd" or "abc"+"d". A protocol that needs messages must
       frame them itself (length prefixes -- Chapter 16's wire format).
     - blocking reads park the process until data arrives, exactly
       like a pipe.

   The ONE socket-specific rule: a socket is FULL-DUPLEX. Both sides
   can read and write at the same time on the same socket -- unlike a
   pipe, which flows one way.
\end{verbatim}

For a server that must hold many connections, blocking reads are the
killer---the exact problem Chapter~10 solved with \code{poll}. The same
medicine works on sockets, because a socket is a descriptor: set it
\emph{non-blocking}, and reads that would block return \code{EAGAIN}
instead. The event loop of Chapter~10 was already the right shape; it
just needed a new source of descriptors.

\section{TCP vs UDP}

Not every conversation wants TCP's guarantees. The alternative,
UDP, trades reliability for speed and simplicity: a UDP socket sends
\emph{datagrams}---individual messages with boundaries, no connection,
no handshake, and no promise that a message arrives, arrives once, or
arrives in order.

\begin{verbatim}
   The two protocols, side by side:

   concern        TCP                              UDP
   ----------     -------------------------------  ---------------------------
   connection     connected, handshake first       connectionless: just send
   delivery       reliable, in order, no dupes     best effort: may lose,
                                                  duplicate, or reorder
   boundaries     byte STREAM: no message edges    DATAGRAM: one message,
                                                  one recvfrom
   overhead       higher (handshake, acks)         minimal
   use for        HTTP, files, databases,          DNS, games, live media,
                  anything needing correctness     broadcasts, fast probes
   API            read() / write()                 recvfrom() / sendto()

   The rule of thumb: if it must be correct, use TCP; if it must be
   fast and small and can tolerate loss, use UDP. Many systems use
   both -- TCP for control, UDP for media.
\end{verbatim}

A UDP server never accepts; it binds and then loops on
\code{recvfrom}, which returns the datagram \emph{and} the address it
came from, so the server can reply with \code{sendto}.

\begin{lstlisting}[caption={udp\_echo.c: a datagram echo server.}]
#include <stdio.h>
#include <stdlib.h>
#include <string.h>
#include <unistd.h>
#include <arpa/inet.h>
#include <sys/socket.h>

#define PORT 9001

int main(void) {
    int fd;
    struct sockaddr_in addr;
    struct sockaddr_in peer;
    socklen_t peer_len;
    char buf[1024];
    ssize_t got;

    fd = socket(AF_INET, SOCK_DGRAM, 0);     /* UDP */
    if (fd < 0) {
        perror("socket");
        return 1;
    }

    memset(&addr, 0, sizeof(addr));
    addr.sin_family = AF_INET;
    addr.sin_port = htons(PORT);
    addr.sin_addr.s_addr = htonl(INADDR_ANY);

    if (bind(fd, (struct sockaddr *) &addr, sizeof(addr)) < 0) {
        perror("bind");
        return 1;
    }

    printf("UDP echo on port %d\n", PORT);

    while (1) {
        peer_len = sizeof(peer);
        got = recvfrom(fd, buf, sizeof(buf), 0,
                       (struct sockaddr *) &peer, &peer_len);
        if (got < 0) {
            perror("recvfrom");
            break;
        }

        /* Each datagram carried its own source: reply to IT. */
        sendto(fd, buf, (size_t) got, 0,
               (struct sockaddr *) &peer, peer_len);
    }

    close(fd);
    return 0;
}
\end{lstlisting}

\begin{verbatim}
   recvfrom / sendto, and why they are different:

     read/write  on a CONNECTED socket: the peer is known, no address
     recvfrom    on an UNCONNECTED socket: each datagram brings its
                 own return address (peer), used by sendto to reply

   One UDP socket can talk to thousands of peers, because each
   datagram is self-contained. That is UDP's superpower and its
   curse: no connection means no one is keeping track of anyone.
\end{verbatim}

\begin{osconcept}
The choice of TCP or UDP is a contract about how much the kernel will
do for you. TCP asks the kernel to run a small protocol machine per
connection---acknowledging, retransmitting, ordering, and closing---
so your application can trust a byte stream. UDP asks the kernel for
almost nothing, and your application inherits the network's
unreliability. The OS provides both and the differences between them
are entirely kernel policy.
\end{osconcept}

\section{The Scalable Server: One Loop, Many Sockets}

The one-client echo server blocks on a single \code{accept}, then on a
single \code{read}---it can serve exactly one client at a time. The
fix is the chapter you already know: make every socket non-blocking,
put them all in a \code{poll} array, and let one loop serve thousands.
This is the event loop of Chapter~10, now watching sockets instead of
pipes, with one new move: when the \emph{listener} reports readiness,
we \code{accept} and add the new client to the set.

\begin{lstlisting}[caption={echo\_server.c: a poll-driven, multi-client echo server.}]
#include <stdio.h>
#include <stdlib.h>
#include <string.h>
#include <unistd.h>
#include <fcntl.h>
#include <poll.h>
#include <errno.h>
#include <arpa/inet.h>
#include <sys/socket.h>

#define PORT 9000
#define MAX_CLIENTS 64

static int set_nonblocking(int fd) {
    int flags;

    flags = fcntl(fd, F_GETFL, 0);
    if (flags < 0) {
        return -1;
    }
    return fcntl(fd, F_SETFL, flags | O_NONBLOCK);
}

int main(void) {
    int listen_fd;
    struct sockaddr_in addr;
    struct pollfd fds[MAX_CLIENTS];
    int nfds;
    char buf[1024];

    listen_fd = socket(AF_INET, SOCK_STREAM, 0);
    if (listen_fd < 0) {
        perror("socket");
        return 1;
    }

    memset(&addr, 0, sizeof(addr));
    addr.sin_family = AF_INET;
    addr.sin_port = htons(PORT);
    addr.sin_addr.s_addr = htonl(INADDR_ANY);

    if (bind(listen_fd, (struct sockaddr *) &addr, sizeof(addr)) < 0) {
        perror("bind");
        return 1;
    }
    if (listen(listen_fd, 16) < 0) {
        perror("listen");
        return 1;
    }
    set_nonblocking(listen_fd);

    printf("echo server on port %d\n", PORT);

    /* fds[0] is the listener; clients are appended after it. */
    fds[0].fd = listen_fd;
    fds[0].events = POLLIN;
    nfds = 1;

    while (1) {
        int r;
        int i;

        r = poll(fds, (nfds_t) nfds, -1);
        if (r < 0) {
            if (errno == EINTR) {
                continue;
            }
            perror("poll");
            break;
        }

        for (i = 0; i < nfds; i++) {
            if ((fds[i].revents & POLLIN) == 0) {
                continue;
            }

            if (fds[i].fd == listen_fd) {
                /* A NEW connection is ready to be accepted. */
                int cfd;
                int slot;

                cfd = accept(listen_fd, NULL, NULL);
                if (cfd < 0) {
                    continue;
                }
                set_nonblocking(cfd);

                slot = nfds;
                if (slot >= MAX_CLIENTS) {
                    close(cfd);        /* out of room: refuse politely */
                    continue;
                }
                fds[slot].fd = cfd;
                fds[slot].events = POLLIN;
                nfds = nfds + 1;
                printf("client %d connected\n", cfd);
            } else {
                /* A client has data, or closed. */
                ssize_t got;

                got = read(fds[i].fd, buf, sizeof(buf));
                if (got <= 0) {
                    /* EOF or error: close it, drop it from the set. */
                    printf("client %d gone\n", fds[i].fd);
                    close(fds[i].fd);
                    fds[i] = fds[nfds - 1];   /* swap-last-into-slot */
                    nfds = nfds - 1;
                    i = i - 1;
                } else {
                    write(fds[i].fd, buf, (size_t) got);   /* echo */
                }
            }
        }
    }

    return 0;
}
\end{lstlisting}

\begin{verbatim}
   The loop's two moves, drawn:

        poll(fds, nfds)  -- one wait for EVERY socket
              |
              +-- listener ready?  accept(); append to fds[]
              +-- client ready?    read(); echo
              +-- client EOF?      close(); remove (swap-last)

   - accept() is called only when the LISTENER is readable, so it
     never blocks.
   - reads happen only when a client is readable, and the sockets are
     non-blocking anyway: no client can stall the world.
   - removing a closed client by swapping the LAST slot into its place
     keeps the array dense in O(1) -- the same trick as open-address
     deletion with tombstones in Chapter 4.
\end{verbatim}

\section{Practical: Growing the Echo Server into a Database}

The echo server is a skeleton; its muscle comes from everything
earlier in the book. The same loop that echoes can \emph{parse} and
\emph{store}, turning the server into a networked database with the
command language of Chapter~14 and the durable storage of Chapter~16.

\begin{verbatim}
   The capstone architecture, all the pieces together:

     client                       server (the poll loop of ch.10/17)
     "SET name ada\n"                 |
        |                             v
        |  TCP ----------------->  read()  -- collect bytes into a
        |                             line buffer (may need SEVERAL
        |                             reads: TCP is a stream!)
        |                             v
        |                     lex() + parse()   (ch.14: the MiniDB
        |                             |            front end)
        |                             v
        |                     db_set()           (ch.16: the dblog,
        |                             |            durable on disk)
        |                             v
        |  <----------------------  write("ok\n")

   Two new lessons the network adds on top of the earlier chapters:
     1. A "line" may arrive split across reads, and several lines may
        arrive in one read. The server must buffer until a full
        record, using a length prefix (ch.16's wire format) to know
        when it has one.
     2. The event loop dispatches every connection's bytes to the
        same parser -- one loop, one parser, many clients, and the
        durable store behind it.
\end{verbatim}

Every capability in that picture was built in this book: descriptors
(Chapter~5), the event loop (Chapter~10), the parser (Chapter~14), the
durable log (Chapter~16). The network did not require a new idea; it
required reusing the old ones across a socket.

\begin{researchnote}[The philosophy of the internet protocols]
The finale deserves the field's founding argument. In 1988, David
Clark, one of the architects of the internet protocols, published
``The Design Philosophy of the DARPA Internet Protocols.'' It is the
clearest statement ever written of \emph{why} the network works the
way it does---and it explains every design choice in this chapter.

\begin{verbatim}
   Clark's first principle, condensed:

       "The internet is designed around a fundamental principle: the
        network itself is deliberately simple and makes no guarantees;
        the intelligence and reliability live at the ENDPOINTS, in
        the hosts, not in the middle of the network."

   In this chapter's terms: the network delivers packets best-effort,
   and TCP--which lives in your kernel, at the endpoint--supplies the
   reliability. The routers in the middle stay dumb and fast; the
   cleverness is at the edges. This is the "end-to-end argument," and
   it is why your socket's kernel does so much work.
\end{verbatim}

Clark's paper is short, opinionated, and remains the best one-page
answer to ``why is the network built this way?'' Reading it closes the
book's arc: from bytes and addresses on one machine, to a reliable
byte stream between two machines, designed by the same principles of
layering and explicit responsibility you have practiced all along.

\paperref{D.~D.~Clark, ``The Design Philosophy of the DARPA Internet
Protocols,'' ACM SIGCOMM 1988}{https://www.cs.princeton.edu/courses/archive/fall06/cos561/papers/clark88.pdf}
\end{researchnote}

\section{Structures in Bytes}

The TCP header you just read is one of the three byte structures of a
network conversation. The other two are the socket address you build
with \code{htons}/\code{htonl}, and the IP header that carries every
segment.

\begin{verbatim}
   struct sockaddr_in (IPv4): 16 bytes
   offset  size  field
   0       2     sin_family     (AF_INET = 2)
   2       2     sin_port       (network byte order!)
   4       4     sin_addr       (network byte order!)
   8       8     padding to the generic sockaddr

   the IPv4 header (no options): 20 bytes
   offset  size  field
   0       1     version + IHL      (4 + 4 bits; 5 words = 20 bytes)
   1       1     DSCP / ECN
   2       2     total length
   4       2     identification
   6       2     flags + fragment offset
   8       1     TTL
   9       1     protocol           (6 = TCP, 17 = UDP)
   10      2     header checksum
   12      4     source address
   16      4     destination address

   the UDP header: 8 bytes
   offset  size  field
   0       2     source port
   2       2     destination port
   4       2     length
   6       2     checksum
\end{verbatim}

A datagram on the wire is IP header + transport header + payload
(20 + 20 + data for TCP, 20 + 8 + data for UDP), every integer in
network byte order. The kernel builds and parses these; your program
only ever sets the two fields of the 16-byte \code{sockaddr\_in},
converted with \code{htons} and \code{htonl}.

\section{Summary}

\begin{verbatim}
   The networking toolbox, complete:

   socket(AF_INET, SOCK_STREAM|SOCK_DGRAM, 0)   create an endpoint
   sockaddr_in + htons/htonl/inet_pton          name it (byte order!)
   bind()      give the socket its address and port
   listen()    tell the kernel to queue connections
   accept()    take one connection -> a NEW socket for that client
   connect()   drive the three-way handshake (client side)
   read/write  the connected socket is a byte stream (ch.5 rules)
   recvfrom/sendto  UDP datagrams with per-message addresses
   poll + non-blocking  one loop, many clients (ch.10, reused)

   Rules that keep network code correct:
     1. put every integer on the wire in network byte order;
     2. a connected socket is a byte stream: buffer until you have a
        full record, using a length prefix;
     3. EOF (read == 0) means the peer closed; POLLHUP, too;
     4. never block on one socket when serving many: non-blocking +
        poll;
     5. accept() only when the listener is ready; it never blocks
        then.
\end{verbatim}

\begin{osconcept}
Networking is the operating system's grand finale abstraction: the
kernel owns the hardware, runs the protocols, holds the handshakes,
and presents every conversation as a descriptor. The end-to-end
principle says the network stays simple and the intelligence lives at
the edges---in the kernels of the two hosts, and in the discipline of
your program. You have now built, byte by byte, process by process,
every layer between a variable in RAM and a message on the wire.
\end{osconcept}

This is where the book's tour ends and your systems programming begins.
You hold all the tools: memory and pointers, structures and algorithms,
processes, signals, pipes, threads, filesystems, durable storage, and
the network. The machine is no longer a mystery; it is a set of
contracts you understand, and the rest is experience.

